\section{Data and method}
\label{sec:data}

This section states what each side of the condition is built from, which components are sourced
and which are swept, and how a reader should weigh what comes out. The distinction between
sourced and swept carries through every number that follows and is the reason the central figure
is labeled scenario rather than measured.

\subsection{Sources}

The labor side of the condition is built from federal taxes over wages and salaries, with the
wage share of adjusted gross income and the realized gains bound that travels with it from the
Statistics of Income \citep{IRSSOI}. The retained wage share is fitted on fourteen biennial
vintages of the worker displacement release \citep{BLSWorkerDisplacement} and solved as a fixed
point, since displacement changes the labor market slack the fit is a function of.

The capital side assembles the components of Definition~\ref{def:tauk}. The entity rate and
the treatment of expensing follow \citet{AcemogluManeraRestrepo2020}; the share of rents booked
in lower-taxed jurisdictions is from \citet{TorslovWierZucman2023}; the share of corporate equity
held in taxable accounts is from \citet{RosenthalAustin2016}, \citet{RosenthalBurke2020} and
\citet{RosenthalMucciolo2024}; the deferral factor between accrual and realization is from
\citet{Gravelle2022}; and the published measured rate the assembly is checked against is
\citet{CBO2014}. The two readings of the rent share are \citet{Barkai2020} and
\citet{KarabarbounisNeiman2014}.

The household side uses the Survey of Consumer Finances \citep{FRBSCF} and, for the distribution
of corporate equity across households, the distributional financial accounts \citep{FRBDFA}. The
claim levels and holder shares the incidence calculation runs on come from the Federal Reserve's
financial accounts \citep{FRBZ1} through the labor backing accounts, which are reported
separately. The bust direction is priced on verified historical aggregates and on the wealth
effect of \citet{ChodorowReichNenovSimsek2021}; the financing mix of the observed investment is
read from the filings of nine registrants, with the off-balance-sheet structures described by
\citet{BIS2026} absent from them. The wage bill and federal receipts used as bases are
\result{WageBillBn} and \result{FedReceiptsBn} billion dollars.

\subsection{Two objects borrowed from the accounts, and one built here}
\label{sec:borrowed}

Two quantities used below are built in a companion paper \companion\ rather than here, and a
reader of this paper alone needs them defined. Each is stated in full there; what follows is
enough to read the results in this one.

\paragraph{The retained wage share, which is built here rather than borrowed.} Written $R$, this
is the share of a displaced wage bill that returns as wages, since not every dollar of displaced
pay leaves the wage tax base. It determines the replacement threshold directly through
Definition~\ref{def:condition}, so it is set out in full rather than cited.

It is the product of two quantities, $R = \rho\,\omega$. The reemployment rate $\rho$ is the share
of displaced workers re-employed, fitted on fourteen biennial vintages of the worker displacement
release \citep{BLSWorkerDisplacement} against labor market slack and solved as a fixed point,
because displacement changes the slack the fit is a function of. The retention factor $\omega$ is
the share of the prior wage kept on re-employment, from \citet{Huckfeldt2022}. At the observed
reemployment rate, $\rho = \result{RhoReemploy}$ and $\omega = \result{OmegaWageKept}$, giving
$R = \result{RetainedWageShare}$.

Three specifications of $\rho$ are carried, the observed rate and two fitted variants, and they
give $R$ between \result{RSensLow} and \result{RSensHigh}. The required rate moves with it,
between \result{RReqSensLow} and \result{RReqSensHigh} percent on the narrower labor tax reading,
and the assembled marginal rate of \result{TauKFederalOnly} percent falls short across the whole of
that range, while the annual rate of Section~\ref{sec:annual} clears it across the whole of it. The four decimal places this parameter is sometimes quoted to conceal how much fitting
sits behind it, which is why the sensitivity is reported here beside the point estimate.

\paragraph{The holder mapping.} The companion paper decomposes the stock of US financial claims
twice over, by the income that services each claim and by who holds, guarantees or owes it. Where
this paper reports which share of a given pool of household debt the federal government stands
behind, that attribution is the companion's holder map rather than a construction of our own. A
claim is assigned to the party that bears the credit risk rather than to the party holding legal
title, so an agency-guaranteed mortgage inside a security owned by a pension fund counts as
federal.

\paragraph{The agency classification.} Whether the housing agencies count as federal is a
classification choice rather than a measurement, and the companion paper reports both readings
because they do not agree. It is the largest single judgement call affecting the incidence result
here, and Section~\ref{sec:limitations} records that it moves the holder attribution by about a
factor of seven. Nothing in this paper's revenue-side result depends on it.

\subsection{Sourced, swept, and why the distinction matters}

Two components of the assembled rate have no single sourced value and are swept across a stated
range instead, the effective state corporate rate and the shareholder-level statutory rate: each
is drawn uniformly over the range given for it in the replication package. A third, the
profit-shifted share, was swept in an earlier version and is now sourced from two independent
estimates, as Section~\ref{sec:levers} sets out; the effect of that change on the lever ranking
is reported there rather than absorbed silently.
Everything built on them is labeled \emph{scenario}, and we report the share of the swept grid
in which the condition closes rather than a probability, because a share of a uniformly weighted
grid is not a probability and reporting it as one would overstate what the sweep can tell anyone.

Two features of the assembly do most of the work and neither is our assumption. Under full
expensing the normal return is relieved at the entity level, so at the margin it bears only the
owner's tax, which is why the debt-financed term is negative throughout the sourced range. And
the rent share has two incompatible readings rather than one uncertain value, so every headline
is reported under both.

\subsection{Status labels, and what was independently rebuilt}
\label{sec:rebuild}

Three labels are used throughout and appear in every table caption. \emph{Rebuilt} means the
quantity was reproduced by an independent computational rebuild: a separate run that had no
access to this paper's code and worked from raw data and a published specification alone, inside
a tolerance fixed before the rebuild began. \emph{Measured} means produced by the pipeline,
clearing every applicable bounds check, but not independently rebuilt. \emph{Scenario} means
arithmetic conditional on an assumption stated where the number appears.

What a rebuild establishes needs saying plainly, because the word invites more than it should. It
establishes computational reproducibility: that the stated inputs and the stated construction
produce the stated number. It says nothing about whether the construction is the right one, or
whether the economic assumptions behind it hold. A quantity can be rebuilt to six decimal places
and still rest on an assumption we have no evidence for, and several in this paper do. The
rebuild of a scenario is a rebuild of arithmetic.

The assembled rate and the retained wage share were both reproduced in full by the third rebuild
round, which matched all \result{RepThreeSealed} quantities it covered,
\result{RepThreeRounding} of them within rounding. The blind was procedural rather than enforced,
so the claim is that these quantities were independently rebuilt under a reported blind protocol,
not that they were independently verified. Appendix~\ref{app:rebuild} gives the round-by-round
detail and the causes of the earlier mismatches.
